\clearpage\fancyhead[L]{\small GAUSS TO BODE.}\hypertarget{native-356-lower}{}\label{native-356-lower}
\begin{center}\textbf{GAUSS TO BODE.}\par\vspace{4pt}\hrule\vspace{4pt}\fontsize{10}{12}\selectfont Astronomisches Jahrbuch für 1815. Page 190\qquad194. Berlin 1812.\par\vspace{4pt}\hrule\vspace{4pt}\end{center}

\begin{flushright}Göttingen, 5 May 1812.\end{flushright}
\par\indent I would like to receive some observations of Pallas from you again this year, as I did last year, around the time of opposition (10 June). That is why I am drawing your particular attention to it and telling you what I have observed so far. I saw the planet again for the first time on 4 April, but did not observe it until the 9th. Yesterday’s observation, like that of 9 April, was unsuitable for determining the declination, so I omit it altogether. The planet is scarcely brighter than the 10th magnitude.

\begin{center}
\fontsize{10}{12}\selectfont\setlength{\tabcolsep}{3pt}\renewcommand{\arraystretch}{1.15}\begin{tabular}{c|c|c|c}
1812 & Mean time & R.A. $\PallasSymbol$ & Northern decl. \\
\hline
April\ \ 9 & $11^{\mathrm u}\;32^{\mathrm m}\;27^{\mathrm s}$ & $169^\circ\;5'\;58''4$ & \\
May\ \ 2 & $10\;36\;6$ & $268\;42\;5.3$ & $21^\circ\;4'\;37''7$ \\
4 & $10\;24\;16$ & $268\;32\;27.4$ & \\
\end{tabular}
\end{center}
\SourceLeaf{357}\par\indent Mr. \textsc{Nicolai} finds the following differences by comparison with the elements on which he based his ephemeris:

\begin{center}
\fontsize{10}{12}\selectfont\setlength{\tabcolsep}{3pt}\renewcommand{\arraystretch}{1.15}\begin{tabular}{c|c|c}
 & R.A. & Decl. \\
\hline
April\ \ 9 & $+30''6$ & \\
May\ \ 2 & $+57.6$ & $-5''4$ \\
4 & $+49.4$ & \\
\end{tabular}
\end{center}

\par\indent I consider my observations good. The comparison was made with 93, 96, 95 in Hercules in \textsc{Piazzi}’s catalogue. I am inclined to attribute the difference in right ascension to the perturbations from Saturn and Mars, which have so far been neglected in my calculations but which I will soon calculate as well. The total number of periodic perturbations due to Jupiter that I have calculated amounts to 400, with all those below $1''$ omitted. Several remarkable results emerged on this occasion, but for now I must hold them back. For my entire calculation can only be regarded as preliminary and must be repeated on a still larger scale.

\par\indent I have yet to observe any of the other new planets. Last night I only brought the region where Juno stands into view in the telescope. It was very low in the sky, and the faint stars were still hard to see. The place where Juno should have stood according to the ephemeris seemed empty, whereas at a right ascension about $25'$ smaller there was a star of the 9th to 10th magnitude, which may have been Juno. This is merely a preliminary note. It will soon be possible to determine the matter more closely.

\par\indent I take this opportunity to share a few more astronomical observations with you. — I learned from Mr. von \textsc{Lindenau} on 9 December last year that the second comet had appeared, and I found it at once that same day. Here are all my observations of it. I consider them all good except for the last declination, which is a little less reliable.

\begin{center}
\fontsize{10}{12}\selectfont\setlength{\tabcolsep}{3pt}\renewcommand{\arraystretch}{1.15}\begin{tabular}{c|c|c|c}
 & Mean time & R.A. & Decl. \\
\hline
1811 Dec.\ \ 9 & $10^{\mathrm u}\;6^{\mathrm m}\;52^{\mathrm s}$ & $63^\circ\;49'\;41''4$ & $10^\circ\;21'\;55''5$ S. \\
11 & $10\;34\;1$ & $63\;33\;18.0$ & $8\;39\;46.4$ \\
12 & $\;8\;5\;52$ & $63\;26\;25.3$ & $7\;54\;25.9$ \\
1812 Jan.\ \ 3 & $\;6\;54\;46$ & $62\;37\;6.4$ & $9\;47\;14.7$ N. \\
4 & $\;6\;54\;15$ & $62\;40\;38.6$ & $10\;29\;57.8$ \\
Feb.\ \ 2 & $11\;51\;52$ & $68\;9\;17.9$ & $26\;1\;34$ \\
\end{tabular}
\end{center}

\Needspace{15\baselineskip}\par\indent Immediately after my observation of 12 December I calculated the parabolic elements. I omit them here, however, and give instead the improved ones that Mr. \textsc{Nicolai}, a very skilful calculator, obtained immediately after my observation of 4 January.

\begin{center}\fontsize{10}{12}\selectfont\setlength{\tabcolsep}{2pt}\renewcommand{\arraystretch}{1.15}\begin{tabular}{lll}
\multicolumn{3}{l}{Time of perihelion Nov 11.39211, Göttingen meridian}\\
— Longitude & $.\;.\;.\;.\;.\;.\;.\;.\;.\;.\;.\;.\;.\;.\;.$ & $47^\circ\;39''\;36''1$\\
— Logarithm of the distance & $.\;.\;.\;.\;.\;.\;.\;.$ & $0.2011781$\\
Longitude of $\AscendingNode$ & $.\;.\;.\;.\;.\;.\;.\;.\;.\;.\;.\;.\;.\;.\;.\;.$ & $92^\circ\;53'\;44''2$\\
Inclination of the orbit & $.\;.\;.\;.\;.\;.\;.\;.\;.\;.\;.\;.\;.\;.$ & $31\;32\;38.7$\\
\multicolumn{3}{l}{Direct motion.}\\
\end{tabular}\end{center}

\Needspace{10\baselineskip}\par\indent According to Mr. \textsc{Nicolai}’s calculation, comparison of these elements with my observations gave the following results.
\SourceLeaf{358}\begin{center}
Difference
\end{center}

\begin{center}
\fontsize{10}{12}\selectfont\setlength{\tabcolsep}{3pt}\renewcommand{\arraystretch}{1.15}\begin{tabular}{c|c|c||c|c|c}
1811 & R.A. & Decl. & 1812 & R.A. & Decl. \\
\hline
Dec. 9 & $-5''0$ & $+3''4$ & Jan. 3 & $-6''5$ & $+0''3$ \\
11 & $+2.9$ & $-32.7$ & 4 & $+2.8$ & $-12.7$ \\
12 & $+8.5$ & $-27.6$ & Feb. 2 & $-113.7$ & $+18.1$ \\
\end{tabular}
\end{center}

\par\indent Mr. \textsc{Nicolai} intends soon to improve his elements once again using the last observations made in February.

\Needspace{30\baselineskip}\par\indent Here are some further observations of stellar occultations:

\begin{center}\fontsize{9}{11}\selectfont\setlength{\tabcolsep}{2pt}\renewcommand{\arraystretch}{1.15}\begin{tabular}{llrrrll}
1811 Oct. 23 & Star of the 6th magnitude & $7^{\mathrm u}$ & $21^{\mathrm m}$ & $40^{\mathrm s}1$ & Ingress & Harding\\
— & 7th & 8 & 14 & 40.4 & — & Harding\\
— & 6th\ 7th & 8 & 30 & 33.8 & — & Gauss\\
24 & 7th\ 8th & 7 & 38 & 48.8 & — & Harding\\
1812 Feb. 1 & $\eta$ Virginis & 12 & 21 & 26.7 & — & Gauss\\
19 & $\gamma$ Tauri & 7 & 22 & 15.0 & Egress & Gerling\\
— & 71 Tauri & 9 & 56 & 3.4 & Ingress & Gerling\\
 &  &  &  & 3.8 & — & Harding\\
 &  &  &  & 3.9 & — & Gauss\\
— & $1\,\theta$ Tauri & 10 & 58 & 56.9 & — & Harding\\
 &  &  &  & 57.4 & — & Gauss\\
 &  &  &  & 57.9 & — & Gerling\\
 &  & 11 & 55 & 34.6 & Egress & Gauss\\
 &  &  &  & 38.6 & — & Harding\\
— & $2\,\theta$ Tauri & 10 & 58 & 12.0 & Ingress & Gauss and Harding\\
 &  & 11 & 59 & 6.6 & Egress & Gauss\\
 &  &  &  & 8.6 & Egress & Harding\\
— & 160 Mayer & 12 & 3 & 39.5 & Ingress & Harding\\
 &  &  &  & 40.0 & — & Gauss\\
— & 162 Mayer & 12 & 26 & 27.0 & — & Gauss\\
 &  &  &  & 27.5 & — & Harding\\
20 & 111 Tauri & 11 & 7 & 46.8 & Ingress & Gerling\\
— & Anonyma & 11 & 36 & 54.0 & — & Harding\\
— & 117 Tauri & 13 & 0 & 25.0 & — & Gerling\\
 &  &  &  & 25.1 & — & Harding\\
\end{tabular}\end{center}

\par\indent The occultation of $\eta\;\text{\AstroSymbols ♍}$ on 23 April did not occur at all here in Göttingen.

\par\indent I still have an observation to report from the occultation of $\lambda\;\text{\AstroSymbols ♒}$ on 2 September 1811. Colonel von \textsc{Tischleder} sent it to me with instructions to pass it on to you.

\par\indent Rosenberg Fortress above Cronach: ingress of $\lambda$ Aquarii, 1811 Sept. 2, $10^{\mathrm u}$ $25^{\mathrm m}$ $5^{\mathrm s}4$ mean time.

\par\indent From this and from the observations at Göttingen and Mannheim, Mr. \textsc{Wachter} calculated the following times of conjunction: Mannheim $10^{\mathrm u}$ $16^{\mathrm m}$ $38^{\mathrm s}01$, Göttingen $10^{\mathrm u}$ $22^{\mathrm m}$ $29^{\mathrm s}69$, Cronach $10^{\mathrm u}$ $28^{\mathrm m}$ $3^{\mathrm s}50$.

\par\indent Thus the longitude of Cronach, from comparison with

\begin{center}
\fontsize{10}{12}\selectfont\setlength{\tabcolsep}{3pt}\renewcommand{\arraystretch}{1.15}\begin{tabular}{ll}
Göttingen & $35^{\mathrm m}$ $55^{\mathrm s}81$ \\
Mannheim & 35 \quad $57.49$ \\
mean & 35 \quad $56.65$ or in degrees $28^\circ$ $59'$ $10''$ \\
\end{tabular}
\end{center}
\SourceLeaf{359}\par\indent The colonel puts the elevation of the pole at $=50^\circ$ $14'$ $42''6$.

\par\indent Work on our new observatory is proceeding briskly this summer too. I expect a 12-inch repeating circle from \textsc{Reichenbach} soon. Dr. \textsc{Gerling} is going to the lyceum in Cassel as a teacher of mathematics. He will take up his post in September. His skill and zeal make it desirable that he should be able to pursue astronomical work there as well. His qualifying dissertation on stellar occultations and eclipses, especially on the solar eclipse of 1820 (7 Sept.), which will be annular here (but not in Berlin), is now being printed.
